% !PRINTED 580
\weberpage{580}
%<<CONT
and can be developed in the following way:
\weq{17}{\frac{\sum\limits_{-\infty,\infty}^{\nu} q^{\left(\nu+\frac12+\frac{h'}{m}\right)^2} e^{\frac{(2\nu+1)\pi i}{m}h}}{\sum\limits_{-\infty,\infty}^{\nu} q^{\left(\nu+\frac{h'}{m}\right)^2} e^{\frac{2\nu\pi i}{m}h}}.}
In the denominator of this expression there occurs only one term with
the lowest power of $q$; one obtains it by putting for $\nu$
the integer that lies nearest to the fraction $-h'/m$.

In the numerator too there occurs in general only one lowest
term, whose coefficient is a unit. In the particular
case where $h'/m$ is an integer, however, two
equal lowest terms occur in the numerator, namely (for $h'=0$) $\nu=0,-1$.
Together these give
\[
q^{\frac14}\,2\cos\frac{h\pi}{m}
\]
and $2\cos\dfrac{h\pi}{m}$ is an algebraic unit. [To be inferred from
Bd. I, § 144, (19).]

Accordingly both the expression (17) and \gesperrt{its
reciprocal value} can be developed in ascending powers of $q$ in such a way
that the coefficient of the first term has the value 1
and all the remaining coefficients are integral algebraic numbers.
% [source sic: the preceding exceptional case gives the generally nontrivial algebraic unit 2 cos(h\pi/m)]
From this one concludes, just as in regard to $A(x)$, that
$B(x)\,C(x)$ can be represented in the following way:
\weq{18}{B(x)\,C(x) = x^{2m^2-2} + b_1 x^{2m^2-4} + \cdots + 1,}
in which the $b_1, b_2 \ldots$ are integral rational functions of $\lambda$ with integral
rational numerical coefficients. Moreover the coefficients
 $b$ that are equally distant from the beginning and from the end are
equal to one another.

Let us consider further the equation arising from (8):
\weq{19}{x A(x) - \sqrt{\varkappa}\,\mathrm{sn}\,mv\,D(x) = 0,}
which in regard to $x$ is of degree $m^2$. Its roots are the
quantities:
\[
\sqrt{\varkappa}\,\mathrm{sn}\,v, \quad \sqrt{\varkappa}\,\mathrm{sn}(v+\Omega_{h,h'}).
\]
The product of these roots is, apart from the sign,
equal to the independent term in this equation, and hence equal to
\[
\pm \sqrt{\varkappa}\,\mathrm{sn}\,mv,
\]
and from this there results:
\weq{20}{\pm\frac{\sqrt{\varkappa}\,\mathrm{sn}\,mv}{\sqrt{\varkappa}\,\mathrm{sn}\,v} = \prod^{h,h'} \sqrt{\varkappa}\,\mathrm{sn}(v+\Omega_{h,h'}).}
